\scanpage{6}
\mspage{5}
Donc la donnée de $(\mathscr F,\leq)$ avec figures de t.f. équivaut : $(M,\leq\deleted{\unsure{X\leq}\mathrel{\tikz[baseline=-.4ex,x=.65em,y=.65em]{\draw (0,1.1)--(1,.5)--(0,-.1) (1,1.1)--(0,.5)--(1,-.1);}}}\text{ (ordre)},\ \mathrel{\tikz[baseline=-.4ex,x=.65em,y=.65em]{\draw (0,1.1)--(1,.5)--(0,-.1) (1,1.1)--(0,.5)--(1,-.1);}}\text{ (\unsure{rel b.|rel\textsuperscript{tion}}, réflexive sym)})$ satisfaisant

a) \deleted{$\unsure{\forall x\in M}$,} \inserted{\unsure{posant}} \deleted{$\widetilde X=\{y\in M\mid y\leq x\}$ \unsure{être}} Toute partie fermée majorée $\Phi$ de $M$ est de type fini (i.e. \deleted{\unsure{$\exists$}}réunion finie d'ens. $\widetilde X=M_{\leq\unsure{x|X}}$)
\margin{On aura alors $\widetilde{\mathscr F}={}$ ens des parties fermées de $M$, de type fini, formées \unsure{de|d'} éléments deux à deux compatibles (\unsure{pour|p. r.} $\mathrel{\tikz[baseline=-.4ex,x=.65em,y=.65em]{\draw (0,1.1)--(1,.5)--(0,-.1) (1,1.1)--(0,.5)--(1,-.1);}}$)}

\[
\text{b)}\quad X\mathrel{\tikz[baseline=-.4ex,x=.65em,y=.65em]{\draw (0,1.1)--(1,.5)--(0,-.1) (1,1.1)--(0,.5)--(1,-.1);}}Y,\ X'\leq X,\ Y'\leq Y\ \Longrightarrow\ X'\mathrel{\tikz[baseline=-.4ex,x=.65em,y=.65em]{\draw (0,1.1)--(1,.5)--(0,-.1) (1,1.1)--(0,.5)--(1,-.1);}}Y'.
\]

% slanted
\margin{Critère de \uline{compatibilité} \uline{par strates}}
\uline{At 3}\quad Soient $F,G\in\mathscr F$ tels que $\forall\deleted{\unsure{x|de}}X\in\widetilde F,\ Y\in\widetilde G$, on ait $X\mathrel{\tikz[baseline=-.4ex,x=.65em,y=.65em]{\draw (0,1.1)--(1,.5)--(0,-.1) (1,1.1)--(0,.5)--(1,-.1);}}Y$, alors $F\mathrel{\tikz[baseline=-.4ex,x=.65em,y=.65em]{\draw (0,1.1)--(1,.5)--(0,-.1) (1,1.1)--(0,.5)--(1,-.1);}}G$ (i.e. $F\vee G$ existe)

\deleted{\unsure{At}} Les axiomes suivants concernent les relations entre $\leq$ et $\ll$. 1

% slanted
\margin{Critère de \uline{raffinement} \uline{par strates}\\ [\uline{Corollaire} $F\leq G\Longrightarrow F\ll G$]}
\uline{At}~4\quad $\forall F,G\in\mathscr F,$
\[
F\ll G\ \Longleftrightarrow\ \forall X\in\widetilde F,\exists Y\in\widetilde G,\ \text{t.q. }X\ll Y
\]
Ramène la connaissance de $\ll$ dans $\mathscr F$ à celle de $\underset{M}{\ll}$

% slanted
\margin{\tikz[x=.6em,y=.6em]{\draw (0,0)..controls(.6,.45)and(1.5,.58)..(2,.35); \draw (2.03,.57)--(2.3,.57);}\enspace Axiome de \unsure{la|le} plus petite strate, ou du \uline{raffinement interne} $\nearrow$}
\uline{At} \deleted{\unsure{5}}6\quad $\deleted{\forall}X,F\in\mathscr F,\ X\triangleleft F$ (i.e $X\in\widetilde F$), soit $F^{X}=\{Y\in\widetilde F\mid X\deleted{\unsure{\leq}}\ll Y\}$\unsure{,|.} Alors $F^{x}$ a un plus petit élément.

Notation
\[
\begin{aligned}
X\overset{\circ}{\ll}Y&\ \overset{\text{\unsure{dfn|df|defn|défn|déf}}}{\Longleftrightarrow}\ X\text{ est plus petit élément de }\widetilde Y^{x}\\
&\text{i.e.}\quad X\ll Y,\quad\text{et } X\ll Y'\triangleleft Y\Longrightarrow Y'=Y\\
(\text{ou } X\in\mathrm{Omb}(Y)^{\circ}) &\\
&\qquad X\ \text{\uline{raffinement interne} de }Y
\end{aligned}
\]

% slanted
\margin{\tikz[x=.6em,y=.6em]{\draw (0,0)..controls(1,.95)and(2.2,1.65)..(3.05,1.9)..controls(3.32,1.99)and(3.55,1.81)..(3.75,1.62); \draw (1.18,2.12)..controls(1.2,1.81)and(1.43,1.43)..(1.61,1.31);}\enspace équivaut à : \inserted{Si $G,H$,} \tikz[x=.6em,y=.6em]{\draw (0,0)--(.25,1.15)..controls(.32,1.45)and(.55,1.45)..(.8,1.48)--(3.3,1.62); \draw (3.3,1.62)--(3.52,1.77)--(3.68,1.58)--(3.43,1.43)--cycle;} \deleted{\unsure{Axiome}} : sont des \unsure{sous-figures|strates} d'une figure $F$, alors $F\ll G\Longleftrightarrow F\leq G$\\ (\uline{Critère du raffinement \unsure{extrinsèq.|intrinsèq.}} de \uline{figures})}
\uline{At} \deleted{\unsure{6}}5\quad Si $X\mathrel{\tikz[baseline=-.4ex,x=.65em,y=.65em]{\draw (0,1.1)--(1,.5)--(0,-.1) (1,1.1)--(0,.5)--(1,-.1);}}Y$ $(X,Y\in M)$ alors $X\ll Y\Longrightarrow X\leq Y$

ou encore, dans un $\widetilde F$, $\leq$ et $\ll$ coïncident,

ou encore, si $F,G\in M$ et $F\mathrel{\tikz[baseline=-.4ex,x=.65em,y=.65em]{\draw (0,1.1)--(1,.5)--(0,-.1) (1,1.1)--(0,.5)--(1,-.1);}}G$, alors $F\ll G\Longleftrightarrow F\leq G$.

% slanted
\margin{\tikz[x=.6em,y=.6em]{\draw (0,0)--(5,-.25); \draw (.45,.12)--(.12,0)--(.4,-.12);}\enspace \uline{Transitivité} du raffinement interne\quad\tikz[x=.6em,y=.6em]{\draw (1.18,1.8)..controls(.04,2)and(-.02,.3)..(.55,.23)..controls(.67,.2)and(.78,.24)..(.79,.32);}}
\uline{At} \deleted{\unsure{6}}7\quad Soient $X,Y,Z\in M$, alors $X\overset{\circ}{\ll}Y\overset{\unsure{\bullet|\circ}}{\ll}Z\Longrightarrow X\overset{\circ}{\ll}Z$

Traduction en termes de $(M,\leq,\ll,\widetilde{\mathscr F})$
\[
\begin{aligned}
\text{\uline{At}}_{M}\,0&\qquad \leq\ \Longrightarrow\ \ll\\
\text{\uline{A}t}_{M}\,1&\qquad \text{Les }\Phi\in\widetilde{\mathscr F}\ \text{fermés pour }\leq\\
\text{\uline{At}}_{M}\,2&\qquad\text{Si }\Phi\in\widetilde{\mathscr F},\ \Psi\ \deleted{\text{\unsure{pa}}}\text{partie fermée de }\Phi\text{ pour }<,\ \text{alors }\Psi\in\widetilde{\mathscr F}\\
&\qquad\text{et }\widetilde{\mathscr F}\neq\phi\quad(\text{i.e. }\phi\in\widetilde{\mathscr F})\\
\text{\uline{At}}_{M}\,3&\qquad\text{Si }\Phi,\Psi\in\widetilde{\mathscr F},\ \text{telles que }\forall X\in\Phi,\forall Y\in\Psi,\ \deleted{\unsure{\{X,Y\}}}X\mathrel{\tikz[baseline=-.4ex,x=.65em,y=.65em]{\draw (0,1.1)--(1,.5)--(0,-.1) (1,1.1)--(0,.5)--(1,-.1);}}Y\ (\text{i.e.}\\
&\qquad\exists\Phi\in\widetilde{\mathscr F},\ X,Y\in\Phi),\ \text{alors }\Phi\mathrel{\tikz[baseline=-.4ex,x=.65em,y=.65em]{\draw (0,1.1)--(1,.5)--(0,-.1) (1,1.1)--(0,.5)--(1,-.1);}}\Psi\quad(\text{i.e. }\Phi\mathbin{\overset{M}{\cup}}\Psi\in\widetilde{\mathscr F}).
\end{aligned}
\]
